%!TEX TS-program = xelatex
%!TEX encoding = UTF-8 Unicode
% Awesome CV LaTeX Template for Cover Letter
%
% This template has been downloaded from:
% https://github.com/posquit0/Awesome-CV
%
% Authors:
% Claud D. Park <posquit0.bj@gmail.com>
% Lars Richter <mail@ayeks.de>
%
% Template license:
% CC BY-SA 4.0 (https://creativecommons.org/licenses/by-sa/4.0/)
%


%-------------------------------------------------------------------------------
% CONFIGURATIONS
%-------------------------------------------------------------------------------
% A4 paper size by default, use 'letterpaper' for US letter
\documentclass[11pt, a4paper]{awesome-cv}

% Configure page margins with geometry
\geometry{left=1.4cm, top=.8cm, right=1.4cm, bottom=1.8cm, footskip=.5cm}

% Color for highlights
% Awesome Colors: awesome-emerald, awesome-skyblue, awesome-red, awesome-pink, awesome-orange
%                 awesome-nephritis, awesome-concrete, awesome-darknight
\colorlet{awesome}{awesome-red}
% Uncomment if you would like to specify your own color
% \definecolor{awesome}{HTML}{CA63A8}

% Colors for text
% Uncomment if you would like to specify your own color
% \definecolor{darktext}{HTML}{414141}
% \definecolor{text}{HTML}{333333}
% \definecolor{graytext}{HTML}{5D5D5D}
% \definecolor{lighttext}{HTML}{999999}
% \definecolor{sectiondivider}{HTML}{5D5D5D}

% Set false if you don't want to highlight section with awesome color
\setbool{acvSectionColorHighlight}{true}

% If you would like to change the social information separator from a pipe (|) to something else
\renewcommand{\acvHeaderSocialSep}{\quad\textbar\quad}


%-------------------------------------------------------------------------------
%	PERSONAL INFORMATION
%	Comment any of the lines below if they are not required
%-------------------------------------------------------------------------------
% Available options: circle|rectangle,edge/noedge,left/right
\name{Wadim}{Müller}
\position{System Software Expert{\enskip\cdotp\enskip}Linux Enthusiast}
\address{An den Mainwiesen 51, 63533 Mainhausen, Germany}

\mobile{(+49) 176-96526090}
\email{wadim.mueller@mail.de}
%\dateofbirth{January 1st, 1970}
%\homepage{www.posquit0.com}
\github{wafgo}
\linkedin{Wadim Mueller}
% \gitlab{gitlab-id}
% \stackoverflow{SO-id}{SO-name}
% \twitter{@twit}
% \x{x-id}
% \skype{skype-id}
% \reddit{reddit-id}
% \medium{madium-id}
% \kaggle{kaggle-id}
% \hackerrank{hackerrank-id}
% \telegram{telegram-username}
% \googlescholar{googlescholar-id}{name-to-display}
%% \firstname and \lastname will be used
% \googlescholar{googlescholar-id}{}
% \extrainfo{extra information}

\quote{“He who stops being better stops being good.” (Oliver Cromwell)}


%-------------------------------------------------------------------------------
%	LETTER INFORMATION
%	All of the below lines must be filled out
%-------------------------------------------------------------------------------
% The company being applied to
\recipient
  {Canonical Recruitment Team}
  {Canonical}
% The date on the letter, default is the date of compilation
\letterdate{\today}
% The title of the letter
\lettertitle{Job Application for Ubuntu Linux Kernel Engineer - Silicon Enablement}
% How the letter is opened
\letteropening {Dear Mr. Grandbois, }
% How the letter is closed
\letterclosing{Sincerely,}
% Any enclosures with the letter
\letterenclosure[Attached]{Curriculum Vitae}


%-------------------------------------------------------------------------------
\begin{document}

% Print the header with above personal information
% Give optional argument to change alignment(C: center, L: left, R: right)
\makecvheader[R]

% Print the footer with 3 arguments(<left>, <center>, <right>)
% Leave any of these blank if they are not needed
\makecvfooter
  {\today}
  {Wadim Müller~~~·~~~Cover Letter}
  {}

% Print the title with above letter information
\makelettertitle

%-------------------------------------------------------------------------------
%	LETTER CONTENT
%-------------------------------------------------------------------------------
\begin{cvletter}

I am writing to express my enthusiastic interest in the Silicon Enablement Ubuntu Linux Kernel Engineer position I found on canonical.com. As an embedded software engineer with over a decade of experience solving technically challenging problems across diverse CPU/SoC platforms, and with direct contributions to the mainline Linux Kernel, I am confident that my skills are an excellent match for the requirements of this role.

In my current position at the startup CMBlu, I am developing a Battery Management System based on Yocto Linux. My responsibilities include writing and integrating Linux Kernel drivers for our custom boards and leading the board bring-up of our mainboard, which is based on the TI AM64xx SoC. This role has further sharpened my skills in low-level driver development and system integration.

Previously, at Continental, I was responsible for integrating SoC-vendor Board-Support-Packages into Greenhills INTEGRITY OS-powered products. A significant achievement was the creation of a full machine emulation for Renesas RCar3 SoC-based boards using QEMU. This required developing emulations for numerous SoC-specific peripherals and the bootrom to successfully boot a proprietary OS from the reset vector, work which can be viewed on my \href{https://github.com/wafgo/qemu/tree/rcar3_port}{QEMU GitHub fork}.

My experience also includes developing a PCI Endpoint and the necessary drivers to pass through block devices between a PCI host and endpoint on the NXP S32G SoC. This involved writing both an endpoint and a PCI device driver for the Linux kernel, and I have posted the patch on the Linux Kernel Mailing List, which can be viewed here: \href{https://lwn.net/Articles/963538/}{lwn.net/Articles/963538/}.

I am also passionate about driving process improvements. When we faced quality issues with our flash-frontend, I proposed and implemented a customized Yocto-based Linux distribution to resolve the problems. To aid in failure analysis for a Classic AUTOSAR project, I developed a unique Call Stack Analysis Tool capable of determining the worst-case stack usage by analyzing machine instructions directly from an ELF image. Its most innovative feature is the ability to detect calls made through function pointers by examining object file relocation entries.

My Master's Thesis further expanded my expertise, involving the development of an FPGA-accelerated Linux system for real-time video feature extraction on a Xilinx Zynq PSoC. This project included designing two IP-Cores in Verilog and developing the corresponding Linux kernel device drivers. The IP-Core designs, kernel drivers, device tree, and demo application are all available on my \href{https://github.com/wafgo}{GitHub profile}.

Thank you for your time and consideration. I am thrilled by the opportunity to bring my unique skills in low-level development and high-level system solutions to Canonical. I am eager to discuss my qualifications further and demonstrate my commitment to developing high-quality software.

\end{cvletter}


%-------------------------------------------------------------------------------
% Print the signature and enclosures with above letter information
\makeletterclosing

\end{document}
